\section{Outils utilises}

% Les consignes demandent une section indiquant quels outils, notamment d'IA
% generative, ont ete utilises et POUR QUELS USAGES.
%
% Rappel des trois conditions posees par l'enseignant :
%   - l'outil doit permettre de faire plus, ou mieux, que sans lui ;
%   - chaque ligne de code et chaque notion doit pouvoir etre expliquee en
%     detail par les membres du groupe ;
%   - ce que produit une IA doit etre verifie et justifie.
%
% A remplir honnetement et precisement : la soutenance verifie naturellement
% ces conditions.

\subsection{Outils de developpement}

Python et NumPy servent aux calculs et aux tableaux ; PyYAML charge les expériences.
Matplotlib produit les figures et la relecture animée, avec Pillow pour le GIF.
pytest vérifie les cas analytiques, les invariants et les interfaces ; Ruff contrôle
le style. GitHub Actions exécute les tests, la reproduction des figures et le contrôle
Docker, puis compile le rapport lorsqu'il change. Les versions sont consignées dans
les archives et les dépendances sont figées pour les environnements de référence.

\subsection{Outils d'intelligence artificielle generative}

Le 21 octobre 2026, Lenny a utilisé Codex pour examiner les PR et les issues,
implémenter le banc de couloir, ses tests et sa relecture, et préparer la documentation
et le suivi des tâches (PR 35 à 37 et documentation associée). La vérification repose
sur la solution analytique d'un agent isolé, le raffinement temporel, les invariants
sur toutes les paires du groupe, la relecture des archives et la CI. Les exports ont
été inspectés visuellement ; les commandes de lecture sont testées sans fenêtre.

Un défaut de résolution du GIF, révélé par la suite complète lorsque les figures
scientifiques modifiaient les réglages Matplotlib, a été corrigé et couvert par un
test de régression. Ces vérifications ne remplacent pas la compréhension du code
par chaque membre : l'explication collective en préparation de soutenance reste à faire.
Les autres usages d'outils et les proportions de contribution restent à renseigner
par les personnes concernées.

La même journée, Codex a aidé Lenny à relire les PR 43 à 54 de Hamza,
reproduire des contre-exemples, proposer des corrections et répondre aux fils de revue.
Les vérifications portent notamment sur les crêtes du guidage eikonal, les anciennes
archives, les sorties définitives, les géométries admissibles et les mesures incomplètes.
Les tests lents de contacts et de densité ont été exécutés séparément de la CI rapide :
l'hypothèse de convergence du jeu de contact a été retirée et la conclusion de densité
non reproduite reste signalée dans la PR 53. La source primaire de Weidmann, les figures
et les PDF du rapport ont également été relus. Une review automatisée et une CI verte
ne valent pas confirmation personnelle de compréhension ni validation collective d'un jalon.
